%% Revision log for hecke-alg-corresp-for-metapletic-SL2-revised.tex
%% Compiles with pdflatex.

\documentclass[12pt,letterpaper]{amsart}
\usepackage[T1]{fontenc}
\usepackage{newtxtext,newtxmath}
\usepackage[letterpaper,margin=1in]{geometry}
\usepackage{amsmath,amssymb}
\usepackage{array,longtable}
\usepackage{microtype}
\usepackage[hidelinks]{hyperref}

\title{Revision log for ``A correspondence for the generalized Hecke algebra
of the metaplectic cover $\overline{SL(2,F)}$, $F$ $p$-adic''}
\author{David Joyner}
\email{wdjoyner@gmail.com}
\date{Log entry 1, for Draft 1, November 11, 2026.}

\begin{document}
\maketitle

\section{Files and provenance}

\begin{itemize}
\item Original: \emph{New York J.\ Math.}\ \textbf{4} (1998), 223--235.
Original source \texttt{howe6.tex}, with an AMS-\LaTeX{} conversion
\texttt{hecke-alg-corresp-for-metapletic-SL2-1998-amsart.tex} and the PDF
\texttt{hecke-alg-corresp-for-metapletic-SL2\_1998.pdf}.
\item Revised paper: \texttt{hecke-alg-corresp-for-metapletic-SL2-revised.tex}
(this log describes Draft~1 of it).
\item Style guide: \texttt{joyner-math-paper-style-guide.md}, Draft~2, in the
top-level folder of the repository.
\item The file was edited in place, at the author's request.
\end{itemize}

\section{AI usage}

Three AI systems took part. None of them is listed as an author.
\begin{enumerate}
\item \textbf{Argon (Google Gemini), reviser.} Argon made the 2026
mathematical corrections to the original text, before the present pass. The
former \verb|\thanks| note of the revised file summarized them as follows:
corrections to the generator relations of Proposition 6.1 (formerly
Proposition 8), explicit level and cocycle splitting conditions for general
$n$ in Proposition 7.2 (formerly Proposition 10), and a clarification of
Savin's Iwahori--Hecke isomorphism in Section 9 (formerly Section 6). This
log does not itemize Argon's changes line by line. A full comparison with the
1998 conversion is still to be done.
\item \textbf{ChatGPT (OpenAI), referee.} ChatGPT refereed the corrected
text. Its report is not in the repository.
\item \textbf{Claude (Anthropic, Claude Opus 5.5), editor.} Claude made the
Draft~1 changes listed in Sections 3 and 4 below. It brought the file into
line with the style guide, and drafted new expository material. It made no
change to any statement, proof, hypothesis or status of a result, except the
rewording of one sentence recorded in Section 3.6. Claude did not referee the
mathematics.
\end{enumerate}
Review by AI systems is not an independent certification of the results.

\section{Changes in Draft 1 (editorial)}

\subsection{Typesetting (guide Section 1)}
\begin{itemize}
\item Class options \texttt{12pt,letterpaper,reqno} (class \texttt{amsart}
kept). Added \texttt{fontenc} (T1), \texttt{newtxtext} and
\texttt{newtxmath} (Times), \texttt{geometry} with letter paper and one-inch
margins, \texttt{microtype}, and \texttt{array}/\texttt{longtable} for the
notation table.
\item Theorems, propositions, lemmas, corollaries, definitions, examples and
remarks now share one counter numbered within sections. Remarks were
previously on their own counter. Equations are numbered within sections.
\item The hand-typed ``As noted in Section 1'' (old Section 4, now Section 7)
became a \verb|\ref| to Section~4.
\item The final compile with \texttt{pdflatex} has no errors, warnings,
undefined references or overfull boxes. The three overfull boxes of the
previous version disappeared with the change to Times fonts and
\texttt{microtype}. No manual line breaks were added.
\end{itemize}

\subsection{Title page, abstract and acknowledgments (guide Sections 5, 6.2,
6.3)}
\begin{itemize}
\item Subtitle added: ``Draft 1. Corrected and expanded version''. The original
title is unchanged.
\item Date line added: ``Original draft: 1998. Corrected and expanded draft:
2026.''
\item The Naval Academy address was removed. The email address
\texttt{wdjoyner@gmail.com} was kept.
\item The \verb|\thanks| note now identifies the original publication and
the date of the revision.
\item The abstract now opens with an AI-disclosure sentence and a pointer to
the Acknowledgments. The rest of the abstract is unchanged.
\item New unnumbered Acknowledgments section. The original thanks to Bill
Wardlaw and the anonymous referee moved there word for word from the
\verb|\thanks| note. The AI credits, the non-certification statement and the
pointer to this log were added.
\item The sentence of the former \verb|\thanks| note describing the 2026
corrections (``This revised version incorporates corrections to\ldots'') was
removed from the paper (guide 6.4). Its content is in Section 2 of this log.
\end{itemize}

\subsection{Citation keys (guide 1.7)}
All 14 keys and bibliography labels were changed, and every \verb|\cite| was
updated (36 citations). The entries themselves are unchanged.

\begin{center}
\begin{tabular}{llll}
Old & New & Old & New\\ \hline
Bo & Bo76 & KP & KP84\\
CR & CR81 & M & Mo88\\
F & Fl80 & RS & RS77\\
FK & FK86 & S & Sa88\\
G & Ge76 & Sch & Sc98\\
H & HM85 & Sil & Si79\\
K & Ku67 & W & Wa80\\
\end{tabular}
\end{center}

Nine references were added for the literature review: BK98, CL25, GGW18,
GS12, Ka21, Pe13, Sa04, TW18 and We18.

\subsection{Structure (guide Section 2)}
\begin{itemize}
\item The unsectioned opening text became Section~1, Introduction, with an
orienting paragraph, the running example and a paragraph on the organization
of the paper.
\item The paragraph of the opening text that surveys earlier work (``The study
of genuine admissible representations\ldots'') moved, word for word, to the
new Section~2.1. Section~2 is a new literature review, split into work up to
1998 and work after 1998.
\item New Section~3, Background, with four subsections.
\item The old Section~1, ``Some notation'', is now Section~4. A table of
notation was added at its start.
\item Old Sections 2--6 are now Sections 5--9.
\item New Section~10, Summary of results, and Section~11, Glossary, followed
by the Acknowledgments and the references.
\end{itemize}

\subsection{Orienting paragraphs (guide 3.2)}
New orienting paragraphs were added at the start of Sections 1, 2, 2.2, 3,
each subsection of 3, and Sections 4, 5, 8, 9, 10 and 11. Sections 6 and 7
already began with an orienting paragraph, and Section 2.1 begins with one
drawn from new text.

\subsection{Wording fix}
In Section~4 the sentence ``In fact, directly from the definition of $\beta$,
for any $g_1,g_2\in B_k$, the Iwahori factorization $B_k=U_k^-T_kU_k^+$
implies that $\overline{B}_k$ splits.'' was not grammatical. At the author's
request it now reads ``In fact, the definition of $\beta$ and the Iwahori
factorization $B_k=U_k^-T_kU_k^+$ imply directly that $\overline{B}_k$
splits.'' The dangling phrase ``for any $g_1,g_2\in B_k$'' was dropped. The
claim is unchanged.

\section{New mathematical and expository text, for refereeing (ground rule 2)}

Every item below is new in Draft~1 and was written by Claude. None of it
changes or adds to the results of the paper. Each item should be refereed.
\begin{enumerate}
\item \textbf{Example 1.1} (running example: $F=\mathbb{Q}_3$, $n=2$, $k=1$).
Values of the quadratic Hilbert symbol, including $(-1,3)=-1$.
\item \textbf{Section 2} (literature review). Summaries of BK98, Sa04, GS12,
TW18, CL25, Ka21, Pe13, We18 and GGW18 are based on their abstracts only. The
claim that the dual group in Weissman's framework depends on the parity of
$n$, in agreement with Section 9, is from memory (verification level 3).
\item \textbf{Section 3.1} and \textbf{Example 3.1}
($\mathcal{H}(F^\times//\mathcal{O}^\times)\cong\mathbb{C}[t,t^{-1}]$).
\item \textbf{Section 3.2} and \textbf{Example 3.2}: properties of the Hilbert
symbol, $\beta(h(a),h(b))=(b,a)$, $\beta(h(-1),h(3))=-1$ and
$\beta(s_1,s_1)=(-1,-1)=1$ in $\mathbb{Q}_3$.
\item \textbf{Section 3.3} and \textbf{Example 3.3}: $[B:B_1]=q^2(q-1)=18$,
$\operatorname{vol}(B_1s_1B_1)=3$, $\operatorname{vol}(B_1s_1s_2B_1)=9$.
\item \textbf{Section 3.4} and \textbf{Example 3.4}: the description of
Borel's produced module as the smooth part of
$\operatorname{Hom}_{\mathcal{H}(H//K)}(C_c(K\backslash H),W)$ is from memory
and should be checked against Bo76. The trivial representation and its
module.
\item \textbf{Notation table} at the start of Section 4.
\item \textbf{Examples 4.2, 4.3, 4.4, 4.6} (for Definition 4.1, trivial
cocycles, genuine functions, Definition 4.5).
\item \textbf{Examples 5.2 and 5.4} (for Lemmas 5.1 and 5.3). Example 5.4 uses
Lemma 6.3(d) to show $f_{s_1}*f_{s_1}\neq f_{-I}$.
\item \textbf{Example 6.4} (relation A(i) in the running example, checked
against the trivial representation).
\item \textbf{Example 7.3} (Proposition 7.2 in the running example).
\item \textbf{Examples 8.3, 8.5, 8.8, 8.11, 8.13} (main theorem, matrix
coefficients, induced modules, supercuspidal and square-integrable
representations, Corollary 8.12). Example 8.11 uses two standard facts stated
without citation: the Steinberg representation of $SL(2,F)$ is
square-integrable, and a representation with a nonzero Iwahori-fixed vector
is a subquotient of an unramified principal series and so is not
supercuspidal.
\item \textbf{Example 9.3}: the Iwahori--Matsumoto quadratic relation
$T_s*T_s=(q-1)T_s+qT_1$ for $PGL(2,\mathbb{Q}_3)$, stated without citation
(verification level 3).
\item \textbf{Section 10} (summary of results and open questions) and
\textbf{Section 11} (glossary).
\end{enumerate}

\section{Concordance of numbering (guide 6.8)}

All \verb|\label|s were kept. The old numbers are those of the previous
compiled version of the revised file, which agree with the 1998 paper.

\begin{center}
\begin{longtable}{lll}
Label & Old & New\\ \hline
(sections) & Section 1 Some notation & Section 4\\
& Section 2 Some lemmas & Section 5\\
\texttt{section:sl2} & Section 3 & Section 6\\
\texttt{section:metaplectic} & Section 4 & Section 7\\
& Section 5 Application & Section 8\\
\texttt{section:savin} & Section 6 & Section 9\\
\texttt{def:1} & Definition 1 & Definition 4.1\\
\texttt{def:2} & Definition 2 & Definition 4.5\\
\texttt{lemma:3} & Lemma 3 & Lemma 5.1\\
\texttt{lemma:general} & Lemma 4 & Lemma 5.3\\
\texttt{lemma:genuine} & Lemma 5 & Lemma 5.5\\
\texttt{lemma:convolve} & Lemma 6 & Lemma 5.6\\
\texttt{lemma:isomorphic} & Lemma 7 & Lemma 5.7\\
\texttt{prop:main} & Proposition 8 & Proposition 6.1\\
\texttt{rem:1} & Remark 1 & Remark 6.2\\
\texttt{lemma:9} & Lemma 9 & Lemma 6.3\\
\texttt{eq:obvious} & equation (1) & equation (6.1)\\
\texttt{rem:torus\_center} & Remark 2 & Remark 7.1\\
\texttt{prop:main2} & Proposition 10 & Proposition 7.2\\
\texttt{thrm:main} & Theorem 11 & Theorem 8.1\\
\texttt{corollary:main} & Corollary 12 & Corollary 8.2\\
\texttt{def:13} & Definition 13 & Definition 8.4\\
\texttt{prop:silberger} & Proposition 14 & Proposition 8.6\\
\texttt{def:15} & Definition 15 & Definition 8.7\\
\texttt{prop:borel} & Proposition 16 & Proposition 8.9\\
\texttt{def:17} & Definition 17 & Definition 8.10\\
\texttt{cor:18} & Corollary 18 & Corollary 8.12\\
\texttt{thm:savin} & Theorem 19 & Theorem 9.1\\
\texttt{cor:20} & Corollary 20 & Corollary 9.2\\
\end{longtable}
\end{center}

\section{Original and new results (guide 6.5)}

Every numbered theorem, proposition, lemma, corollary, definition and remark
is from the original 1998 paper, as corrected by Argon in 2026 (Section 2).
Draft~1 establishes no new results. The examples, the background, the
literature review, the summary and the glossary are new exposition (Section
4).

\section{Sources not consulted (guide 4.3)}

Claude consulted none of the 14 original references, so no citation in the
paper was checked against its source. In particular, these precise citations
are unchecked: Ku67 \S1 and Theorem~2, HM85 Chapter~3 (Proposition~2.2,
Theorem~2.1, page~33 and pages~38--39), Mo88 Proposition~3.4, Si79
Corollary~1.10.6, CR81 Lemma~3.14, and Bo76 \S\S1--2.

For the nine new references, the bibliographic data (authors, title,
journal, volume, year and pages) were checked against the publisher's or
arXiv's web page, except BK98 and Sa04, which were taken from the reference
list of GS12 as shown on the AMS site. The summaries rely on the abstracts
only. The full texts were not read.

\section{Points for the author and the referees}

These were noticed during the editorial pass. None was changed, because each
would need a mathematical change (ground rule 1).
\begin{enumerate}
\item Lemma 5.7 assumes only $s(1)=1$. The paragraph before it also requires
$s|_{B_k}\equiv1$, which seems to be needed for $\phi$ to preserve
bi-$B_k^*$-invariance.
\item Proofs that are omitted or only sketched (guide 3.7): C(i) in the proof
of Proposition 7.2 (``is omitted''), C(ii) and C(iii) there (sketched), the
proofs of Propositions 8.6 and 8.9 (``one need only check''), the second part
of Corollary 8.12 (``omitted''), and Corollary 9.2, which has no proof.
\item Remark 6.2(3) (``Possibly B(iii) is a superfluous relation'') is a
speculation. It is listed as an open question in Section 10.
\item The symbols $U_k^\pm$ and $T_k$ are used in Section 4 before Lemma 6.3
and Remark 7.1 define them (guide 3.5). The new notation table now lists
them before their first use. The text itself was not changed.
\end{enumerate}

\section{Style-guide rules: defaults and exceptions}

\textbf{[DECIDE] defaults used.} Document class: kept \texttt{amsart}.
Times package: \texttt{newtx}. Glossary placement: notation near the
beginning (Section 4), glossary at the end just before the acknowledgments
(Section 11). AI systems on the author line: not listed, with the credit in
the acknowledgments.

\textbf{Rules not applied, and why.}
\begin{itemize}
\item The paper-specific rules in guide Section 9 for the metaplectic
Jacquet--Langlands series (status labels (P), (E), (C) and so on, hypothesis
tags, the $|\cdot|_p$ convention, a new file for each draft). This is a
standalone 1998 paper stored in that repository. It is not part of the
series. At the author's request, the file was edited in place.
\item Rule 6.6 (``Remark 2026.$n$''). Choosing what to mark would require a
verification that this editorial pass did not do. The candidates are listed
in Section 8, item 2.
\item Rules 3.6 and 3.7 (state all hypotheses, do not skip steps). Meeting them
would change statements or proofs. See Section 8.
\item Rule 3.3 (an example in every section) was not applied to the literature
review, the summary of results or the glossary, which explain no new concepts.
\item Rule 3.8 (no semicolons) was applied to new text only. Existing
semicolons were kept.
\item Rule 8.6 (copyright). The repository is public and already contains the
corrected version of the published paper. The copyright position has not
been checked.
\end{itemize}

\end{document}
